\chapter{Validation, Testing, and Results}

\section{Introduction}
This chapter presents the validation methodology and the main technical results of \projectname{}. Validation was designed to determine whether the implemented platform was operational across its complete execution path: service readiness, orchestration, distributed storage and processing, governed transformation, data-quality control, analytical access, observability, and secured interaction. The assessment therefore relies on retained runtime evidence rather than architectural intent alone \cite{kleppmann2017ddia,apache_airflow_docs,prometheus_docs}.

The quantitative values reported in this chapter correspond to the retained validation capture of August 18, 2026. They represent the observed state of that execution cycle rather than permanent capacity or performance guarantees.

\section{Validation Strategy}
The validation campaign covered six complementary axes, summarized in Table~\ref{tab:chapter8-validation-strategy}.

The strategy followed a gated progression. Infrastructure readiness was assessed first because later results are not reliable when storage, metadata, compute, query, identity, or monitoring services are unavailable. The five Airflow workflows were then checked in their implemented order. Multi-node participation was evaluated separately from workflow completion so that a successful task could not be mistaken for proof of distributed execution.

Acceptance combined Airflow histories for execution state, service interfaces for cluster roles, persisted artifacts for durable outputs, and external Trino queries for analytical usability. Dashboard values were treated as time-bound observations, and conclusions were limited to capabilities documented in Chapter~IV and Appendix~\ref{app:validation-evidence}.

\begin{table}[!htbp]
    \centering
    \footnotesize
    \renewcommand{\arraystretch}{1.16}
    \arrayrulecolor{tablelineblue}
    \rowcolors{2}{tablebodyblue}{white}
    \begin{tabularx}{\textwidth}{|Y|Y|Y|}
        \hline
        \rowcolor{tableheaderblue}
        \textcolor{white}{\textbf{Validation axis}} & \textcolor{white}{\textbf{What was checked}} & \textcolor{white}{\textbf{Acceptance evidence}} \\
        \hline
        Infrastructure & service endpoints, cluster membership, readiness, and storage state & reachable services and healthy distributed members \\
        \hline
        Orchestration & five-DAG execution order, task completion, and rerun behavior & successful workflow histories and controlled lifecycle transitions \\
        \hline
        Distributed execution & Spark workers, HDFS DataNodes, Kafka brokers, and Trino roles & runtime participation across separate virtual machines \\
        \hline
        Transformation and quality & dbt models, tests, Great Expectations checkpoints, and retained artifacts & curated tables, validation JSON, and HTML Data Docs \\
        \hline
        Security & Kerberos, SPNEGO/HTTPS, gateways, SSH, and identity boundaries & functioning protected access paths with documented scope limits \\
        \hline
        Observability and access & metrics, dashboards, run history, and external SQL querying & measurable service state and successful analytical consumption \\
        \hline
    \end{tabularx}
    \caption{Validation strategy adopted for the deployed platform}
    \label{tab:chapter8-validation-strategy}
\end{table}
\FloatBarrier

\section{Infrastructure Validation}
Infrastructure validation established the operational baseline before pipeline results were interpreted. The campaign covered Kafka, HDFS, Hive Metastore, Spark, Trino, Kerberos, the browser-facing gateways, and the monitoring stack \cite{apache_kafka_docs,apache_hadoop_docs,apache_hive_docs,apache_spark_docs,trino_docs,mit_kerberos_docs,prometheus_docs,grafana_docs}.

The retained checks confirmed that:
\begin{itemize}
    \item Kafka brokers were reachable on VM1, VM2, and VM3;
    \item HDFS exposed an active NameNode and three live DataNodes;
    \item Hive Metastore was available as the shared metadata authority;
    \item the Spark master, three workers, Thrift Server, and History Server were reachable;
    \item the Trino coordinator and its two dedicated workers were available through the protected entry path;
    \item Kerberos and the secured HDFS endpoints were operational;
    \item Prometheus, Grafana, and the configured exporters were reachable from the control plane.
\end{itemize}

The HDFS namespace contained two principal branches: \texttt{/bronze} for replayable raw ingestion and \texttt{/warehouse} for governed analytical storage. The client-specific bronze branch contained the three retained datasets, while the Iceberg warehouse contained the raw, staging, intermediate, and analytics namespaces.

\begin{table}[!htbp]
    \centering
    \footnotesize
    \renewcommand{\arraystretch}{1.14}
    \arrayrulecolor{tablelineblue}
    \rowcolors{2}{tablebodyblue}{white}
    \begin{tabularx}{\textwidth}{|L{0.39\textwidth}|Y|}
        \hline
        \rowcolor{tableheaderblue}
        \textcolor{white}{\textbf{Observed HDFS indicator}} & \textcolor{white}{\textbf{Retained runtime result}} \\
        \hline
        DataNode state & 3 live nodes, 0 dead nodes \\
        \hline
        Configured storage capacity & 62.97 GB \\
        \hline
        DFS used & 918.83 MB (1.42\%) \\
        \hline
        DFS remaining & 7.22 GB (11.46\%) \\
        \hline
        Filesystem objects & 338 objects: 196 files/directories and 142 replicated blocks \\
        \hline
        DataNode usage distribution & minimum 1.19\%, median 1.48\%, maximum 1.61\%, standard deviation 0.18\% \\
        \hline
        Replication health & 0 under-replicated blocks and 0 pending deletions \\
        \hline
        Storage type & DISK across 3 nodes in service \\
        \hline
    \end{tabularx}
    \caption{HDFS runtime indicators retained from the distributed deployment}
    \label{tab:chapter8-hdfs-runtime-evidence}
\end{table}

Configured capacity also includes host storage classified outside DFS. Consequently, DFS used and DFS remaining should not be interpreted as a direct arithmetic decomposition of the configured total.

\begin{table}[!htbp]
    \centering
    \footnotesize
    \renewcommand{\arraystretch}{1.14}
    \arrayrulecolor{tablelineblue}
    \rowcolors{2}{tablebodyblue}{white}
    \begin{tabularx}{\textwidth}{|L{0.29\textwidth}|L{0.31\textwidth}|Y|}
        \hline
        \rowcolor{tableheaderblue}
        \textcolor{white}{\textbf{Observed path}} & \textcolor{white}{\textbf{Observed objects}} & \textcolor{white}{\textbf{Validation meaning}} \\
        \hline
        \texttt{/} & \texttt{/bronze}, \texttt{/warehouse} & separates replayable persistence from governed analytical storage \\
        \hline
        \path{/bronze/client_1} & \path{bank_marketing}, \path{online_retail_2}, \path{online_shoppers_intention} & confirms materialized bronze data for the three retained sources \\
        \hline
        \texttt{/warehouse/iceberg} & \texttt{raw\_data.db}, \texttt{staging.db}, \texttt{intermediate.db}, \texttt{analytics.db} & confirms the physical lakehouse-layer progression \\
        \hline
    \end{tabularx}
    \caption{HDFS namespace evidence retained from the validated deployment state}
    \label{tab:chapter8-hdfs-namespace-evidence}
\end{table}

Together, the service and namespace observations establish that the storage layer was active, populated, distributed, and free from immediate replication anomalies during the retained capture.
\FloatBarrier

\section{Pipeline, Transformation, and Quality Validation}
Pipeline validation followed the five prioritized Airflow workflows implemented for the final platform:
\begin{enumerate}
    \item \path{customerdna_client1_environment_reset_pipeline};
    \item \path{customerdna_client1_lakehouse_setup_pipeline};
    \item \path{customerdna_client1_lakehouse_readiness_pipeline};
    \item \path{customerdna_client1_kafka_hdfs_spark_lakehouse_pipeline};
    \item \path{customerdna_client1_dbt_spark_lakehouse_pipeline}.
\end{enumerate}

This decomposition validates lifecycle control without relying on one monolithic workflow. Reset returns the platform to a controlled baseline, setup initializes the required storage and lakehouse structures, readiness prevents deeper execution against unavailable services, the Kafka--HDFS--Spark workflow materializes the distributed raw foundation and executes its embedded Great Expectations tasks, and the dbt workflow builds and tests the governed models \cite{apache_airflow_docs}.

\begin{table}[!htbp]
    \centering
    \scriptsize
    \renewcommand{\arraystretch}{1.12}
    \arrayrulecolor{tablelineblue}
    \rowcolors{2}{tablebodyblue}{white}
    \begin{tabularx}{\textwidth}{|L{0.34\textwidth}|L{0.13\textwidth}|L{0.15\textwidth}|Y|}
        \hline
        \rowcolor{tableheaderblue}
        \textcolor{white}{\textbf{Validated workflow}} & \textcolor{white}{\textbf{Run date}} & \textcolor{white}{\textbf{Observed duration}} & \textcolor{white}{\textbf{Retained completion evidence}} \\
        \hline
        \path{customerdna_client1_environment_reset_pipeline} & 2026-08-18 14:09:30 & 00:05:02 & successful reset of Kafka topics, bronze storage, and Iceberg lakehouse state; validation tasks for Hive Metastore and Trino also succeeded \\
        \hline
        \path{customerdna_client1_lakehouse_setup_pipeline} & 2026-08-18 14:15:32 & 00:00:15 & successful creation and verification of the bronze namespace and Iceberg warehouse namespaces \\
        \hline
        \path{customerdna_client1_lakehouse_readiness_pipeline} & 2026-08-18 14:16:13 & 00:00:08 & successful readiness checks for HDFS, Hive Metastore, Spark cluster, Spark Thrift, and Trino \\
        \hline
        \path{customerdna_client1_kafka_hdfs_spark_lakehouse_pipeline} & 2026-08-18 14:16:50 & 00:20:15 & successful raw loading into HDFS bronze and Iceberg raw tables, followed by Great Expectations bootstrap and validation tasks \\
        \hline
        \path{customerdna_client1_dbt_spark_lakehouse_pipeline} & 2026-08-18 14:41:17 & 00:14:59 & successful dbt model execution, dbt tests, Spark Thrift validation, and final Trino query validation \\
        \hline
    \end{tabularx}
    \caption{Retained Airflow workflow outcomes for the validated execution cycle}
    \label{tab:chapter8-airflow-run-evidence}
\end{table}

The retained Airflow history demonstrates repeatable, dependency-aware execution rather than a one-time manual sequence. Each lifecycle stage can be diagnosed and rerun at its workflow boundary, while the raw-quality checks remain separately inspectable as named tasks within the fourth workflow. The captured task timings also show that the most expensive steps were the Spark-based raw materialization stage (11 minutes 59 seconds inside the fourth workflow) and the dbt build stage (13 minutes 31 seconds inside the fifth workflow), which is consistent with the fact that these phases perform the heaviest storage and transformation work.

\subsection*{Transformation and Data-Quality Evidence}
Transformation validation confirmed that the raw Iceberg tables could be converted into the implemented staging, intermediate, and analytics model layers through dbt-spark. The retained execution traces included model builds and dbt tests, while Trino inspection confirmed that the resulting governed tables were queryable.

Great Expectations validated the three raw-lakehouse source families through the active checkpoint definitions shown in Table~\ref{tab:chapter8-quality-evidence}. The fourth DAG exposes two explicit quality-task boundaries:
\begin{itemize}[leftmargin=*,itemsep=0.1em,topsep=0.2em]
    \item \path{bootstrap_raw_lakehouse_quality_assets};
    \item \path{validate_raw_lakehouse_quality}.
\end{itemize}
Each execution produced machine-readable validation JSON and HTML Data Docs, separating quality evidence from Airflow task status alone \cite{great_expectations_docs,dbt_docs}.

\begin{table}[!htbp]
    \centering
    \scriptsize
    \renewcommand{\arraystretch}{1.13}
    \arrayrulecolor{tablelineblue}
    \rowcolors{2}{tablebodyblue}{white}
    \begin{tabularx}{\textwidth}{|L{0.20\textwidth}|L{0.52\textwidth}|Y|}
        \hline
        \rowcolor{tableheaderblue}
        \textcolor{white}{\textbf{Source family}} & \textcolor{white}{\textbf{Active checkpoint definition}} & \textcolor{white}{\textbf{Retained evidence}} \\
        \hline
        Bank Marketing & \path{raw_lakehouse_quality_checkpoint__bank_marketing.json} & validation JSON and Data Docs \\
        \hline
        Online Shoppers Intention & \path{raw_lakehouse_quality_checkpoint__online_shoppers_intention.json} & validation JSON and Data Docs \\
        \hline
        Online Retail II & \path{raw_lakehouse_quality_checkpoint__online_retail_2.json} & validation JSON and Data Docs \\
        \hline
    \end{tabularx}
    \caption{Source-specific Great Expectations evidence retained by the integrated raw-pipeline validation tasks}
    \label{tab:chapter8-quality-evidence}
\end{table}

The retained dashboard state also reported three refreshed raw outputs with no skipped or partial incremental outputs. Across the three datasets, approximately 1.12 million raw records were visible: about 41.2K bank-marketing rows, 1.07M online-retail rows, and 12.3K online-shopping-session rows.

\begin{table}[!htbp]
    \centering
    \footnotesize
    \renewcommand{\arraystretch}{1.13}
    \arrayrulecolor{tablelineblue}
    \rowcolors{2}{tablebodyblue}{white}
    \begin{tabularx}{\textwidth}{|L{0.27\textwidth}|L{0.18\textwidth}|Y|}
        \hline
        \rowcolor{tableheaderblue}
        \textcolor{white}{\textbf{Retained raw table}} & \textcolor{white}{\textbf{Observed row volume}} & \textcolor{white}{\textbf{Interpretation}} \\
        \hline
        \texttt{bank\_marketing} & about 41.2K rows & confirms successful landing and raw-lakehouse persistence for the campaign-response source family \\
        \hline
        \texttt{online\_retail\_2} & about 1.07M rows & confirms the platform handled the largest retained source volume within the validated run set \\
        \hline
        \texttt{online\_shoppers\_intention} & about 12.3K rows & confirms successful loading for the browsing-intention source family \\
        \hline
    \end{tabularx}
    \caption{Observed raw-lakehouse row volumes retained by the monitoring layer}
    \label{tab:chapter8-raw-row-volumes}
\end{table}
\FloatBarrier

\section{Distributed Execution Validation}
Distributed execution was assessed through the runtime roles of Spark, HDFS, Kafka, and Trino. The Spark master registered workers on VM1, VM2, and VM3; HDFS stored blocks through DataNodes on all three machines; Kafka brokers were distributed across the same nodes; and Trino separated coordination on VM2 from query workers on VM1 and VM3.

The retained Spark capture recorded three live workers, three total cores, 3.0~GiB of aggregate worker memory, an active Thrift service, and twelve completed applications associated with client-data loading and serving. The same retained view also showed the \textit{Thrift JDBC/ODBC Server} as the active long-running service application while the bronze-to-Iceberg jobs for the three retained source families appeared in the completed-application history. Workload allocation was not identical across every worker, which is expected for light jobs; the relevant result is that executor capacity existed beyond the coordinator node \cite{apache_spark_docs}.

\begin{table}[!htbp]
    \centering
    \footnotesize
    \renewcommand{\arraystretch}{1.14}
    \arrayrulecolor{tablelineblue}
    \rowcolors{2}{tablebodyblue}{white}
    \begin{tabularx}{\textwidth}{|L{0.37\textwidth}|Y|}
        \hline
        \rowcolor{tableheaderblue}
        \textcolor{white}{\textbf{Distributed indicator}} & \textcolor{white}{\textbf{Observed runtime result}} \\
        \hline
        Registered Spark workers & 3 alive workers on VM1, VM2, and VM3; 0 dead and 0 decommissioned \\
        \hline
        Spark capacity at capture time & 3 total cores, 1 active core; 3.0 GiB total worker memory, 512 MiB used \\
        \hline
        Spark SQL service & Thrift JDBC/ODBC Server active on the cluster \\
        \hline
        Completed Spark applications & 12 completed applications visible in the retained master/history views \\
        \hline
        Raw-to-Iceberg jobs & successful applications for \texttt{bank\_marketing}, \texttt{online\_retail\_2}, and \texttt{online\_shoppers\_intention} \\
        \hline
        Running service application & \texttt{Thrift JDBC/ODBC Server} active as the SQL entry point for Spark-backed model execution \\
        \hline
        HDFS participation & 3 live DataNodes with replicated blocks across the cluster \\
        \hline
        Kafka participation & one broker on each of VM1, VM2, and VM3 \\
        \hline
        Trino participation & coordinator on VM2 and two dedicated workers on VM1 and VM3 \\
        \hline
    \end{tabularx}
    \caption{Runtime indicators of multi-node storage, transport, processing, and querying}
    \label{tab:chapter8-distributed-execution-evidence}
\end{table}
\FloatBarrier

\section{Security Validation}
Security validation checked whether the hardened paths described in Chapter VII were usable in the deployed environment \cite{mit_kerberos_docs,apache_hadoop_docs,trino_docs,openssh_docs,cryptsetup_docs}.

\begin{table}[!htbp]
    \centering
    \footnotesize
    \renewcommand{\arraystretch}{1.14}
    \arrayrulecolor{tablelineblue}
    \rowcolors{2}{tablebodyblue}{white}
    \begin{tabularx}{\textwidth}{|L{0.29\textwidth}|L{0.30\textwidth}|Y|}
        \hline
        \rowcolor{tableheaderblue}
        \textcolor{white}{\textbf{Security path}} & \textcolor{white}{\textbf{Validated result}} & \textcolor{white}{\textbf{Scope interpretation}} \\
        \hline
        Kerberos-secured Hadoop & service principals, keytabs, and authenticated Hadoop interactions were operational & implemented service-identity baseline \\
        \hline
        HDFS browser access & Windows Kerberos ticket and Firefox SPNEGO reached HTTPS-only HDFS interfaces & client-sensitive path requiring documented hostnames and browser configuration \\
        \hline
        TLS browser exposure & HDFS, Trino, and control-plane interfaces used protected browser entry paths & internal/self-signed trust rather than enterprise PKI \\
        \hline
        Trino analytical access & authenticated gateway allowed controlled access to the coordinator path & gateway authentication, not full enterprise IAM \\
        \hline
        Spark submission & Airflow reached the Spark submit client on VM2 through SSH & authenticated control-plane/data-plane bridge \\
        \hline
        Operational identities & service, administration, and demonstration identities were separated & baseline role separation without centralized secret management \\
        \hline
        Runtime storage & LUKS-backed deployment support was available for VM runtime directories & optional data-at-rest hardening when enabled \\
        \hline
    \end{tabularx}
    \caption{Validated security paths and their implementation boundaries}
    \label{tab:chapter8-security-validation}
\end{table}

The Windows/Firefox SPNEGO workflow depends on the client ticket cache, hostname aliases, and browser negotiation settings. Client-side integration failures must therefore be distinguished from the health of the Kerberos-secured Hadoop back end.
\FloatBarrier

\section{Observability and Analytical Access Validation}
Observability validation confirmed that platform diagnosis was not limited to manual container-log inspection. Prometheus collected infrastructure, exporter, container, and pipeline metrics; Grafana exposed the retained dashboards; cAdvisor supplied node/container observations from all three VMs; and Airflow retained task histories and logs \cite{prometheus_docs,grafana_docs,apache_airflow_docs}.

\begin{table}[!htbp]
    \centering
    \footnotesize
    \renewcommand{\arraystretch}{1.14}
    \arrayrulecolor{tablelineblue}
    \rowcolors{2}{tablebodyblue}{white}
    \begin{tabularx}{\textwidth}{|L{0.39\textwidth}|Y|}
        \hline
        \rowcolor{tableheaderblue}
        \textcolor{white}{\textbf{Monitoring indicator}} & \textcolor{white}{\textbf{Observed metric or state}} \\
        \hline
        Prometheus target health & 6 healthy targets \\
        \hline
        Observed host capacity & 12 CPU cores and 11.46 GiB memory \\
        \hline
        Hive Metastore PostgreSQL & available, 11.4 MiB database size, 44 active connections at capture time \\
        \hline
        Pipeline status & retained Grafana pipeline panels reported \textit{Success} for lakehouse setup, lakehouse readiness, the Kafka--HDFS--Spark raw pipeline, and the dbt-spark transformation workflow \\
        \hline
        Latest visible raw volume & approximately 1.12 million records across the three retained datasets \\
        \hline
        Raw load outcomes & 3 refreshed outputs, 0 skipped, and 0 partial incremental outputs \\
        \hline
        Great Expectations evidence & successful source-specific checkpoint artifacts retained for the raw lakehouse across all three active source families \\
        \hline
    \end{tabularx}
    \caption{Observed monitoring evidence supporting operational visibility}
    \label{tab:chapter8-monitoring-evidence}
\end{table}

Analytical access was validated independently from the transformation engines through an external SQL client. The retained query targeted the Online Retail II raw table through Trino and completed successfully. This matters because it proves that the lakehouse state exposed by the platform was not only internally visible to Spark and dbt, but also consumable from an external analyst-facing tool.

\begin{table}[!htbp]
    \centering
    \footnotesize
    \renewcommand{\arraystretch}{1.14}
    \arrayrulecolor{tablelineblue}
    \rowcolors{2}{tablebodyblue}{white}
    \begin{tabularx}{\textwidth}{|L{0.36\textwidth}|Y|}
        \hline
        \rowcolor{tableheaderblue}
        \textcolor{white}{\textbf{Analytical-access indicator}} & \textcolor{white}{\textbf{Observed runtime result}} \\
        \hline
        Trino version and environment & version 435, environment \texttt{DISTRIBUTED\_DEMO} \\
        \hline
        Coordinator uptime & 6.74 days at capture time \\
        \hline
        Distributed query roles & coordinator on VM2; dedicated workers on VM1 and VM3 \\
        \hline
        Query state & external SQL query completed with state \textit{Finished} \\
        \hline
        Query throughput snapshot & 89.0 rows/s and 242K bytes/s \\
        \hline
        Query target & \texttt{lakehouse.raw\_data.online\_retail\_2}, limited to 200 rows \\
        \hline
        Consumption result & governed lakehouse tables were browsable from an external SQL client \\
        \hline
    \end{tabularx}
    \caption{Runtime evidence that the distributed analytical-access layer was usable}
    \label{tab:chapter8-trino-query-evidence}
\end{table}
\FloatBarrier

\section{Main Results}
The validation campaign confirmed the delivery of a distributed, orchestrated, governed, observable, and security-aware customer-data lakehouse platform under the \projectname{} scope.

The principal result is the demonstrated continuity of the implemented data path. The retained datasets were transported through Kafka, persisted in the HDFS bronze zone, materialized as governed Iceberg tables, transformed through the dbt model layers, and consumed externally through Trino. Consequently, the analytical-access result was not produced by an isolated demonstration dataset; it depended on the same storage, metadata, processing, transformation, validation, and orchestration chain validated throughout the chapter.

Completion was also established through complementary forms of evidence. Airflow histories confirmed workflow order and task state, cluster interfaces exposed the participating distributed services, persisted files and tables demonstrated durable outputs, Great Expectations retained independently inspectable validation artifacts, and Trino supplied a consumer-side query result. This triangulation reduces the risk of interpreting a successful task status or an available endpoint as sufficient proof of end-to-end platform behavior.

\begin{table}[!htbp]
    \centering
    \footnotesize
    \renewcommand{\arraystretch}{1.14}
    \arrayrulecolor{tablelineblue}
    \rowcolors{2}{tablebodyblue}{white}
    \begin{tabularx}{\textwidth}{|L{0.31\textwidth}|Y|}
        \hline
        \rowcolor{tableheaderblue}
        \textcolor{white}{\textbf{Result dimension}} & \textcolor{white}{\textbf{Observed completion evidence}} \\
        \hline
        Orchestration lifecycle & the five workflows covered reset, setup, readiness, integrated raw loading and quality validation, and governed transformation \\
        \hline
        Distributed storage & 3 live DataNodes, populated bronze and warehouse namespaces, 62.97 GB configured capacity, and no immediate replication anomalies \\
        \hline
        Distributed compute & the Spark master registered workers on all three VMs and retained 12 completed applications \\
        \hline
        Governed transformation & dbt materialized staging, intermediate, and analytics models and executed associated tests \\
        \hline
        Data quality & three source-specific Great Expectations checkpoints produced validation JSON and HTML Data Docs \\
        \hline
        Analytical access & Trino served a completed external SQL query through its coordinator and two dedicated workers \\
        \hline
        Observability & Prometheus and Grafana exposed 6 healthy targets, 12 CPU cores, 11.46 GiB observed host memory, metastore state, pipeline status, and quality indicators \\
        \hline
        Data volume & approximately 1.12 million raw records were visible across the three retained datasets \\
        \hline
        Security posture & Kerberos, SPNEGO/HTTPS, authenticated gateways, SSH submission, role separation, and optional LUKS support hardened the deployed baseline \\
        \hline
    \end{tabularx}
    \caption{Consolidated validation results for the implemented platform}
    \label{tab:chapter8-final-results-summary}
\end{table}
\FloatBarrier

The validated platform is a \textbf{Customer 360-ready analytical foundation}, not a finished enterprise AI product or customer master. The retained results show that heterogeneous data could enter through Kafka, persist in distributed bronze storage, become governed lakehouse tables, pass modeled transformations and quality controls, remain queryable through Trino, and expose operational evidence through the monitoring stack. These outcomes establish the completed Data Engineering and Big Data contribution while preserving the boundary between the implemented platform and future AdOptimizer AI capabilities.

\section{Conclusion}
This chapter confirmed that \projectname{} was tested as an integrated operational platform rather than as a collection of independently reachable services. Readiness, orchestration, multi-node participation, transformation, quality, analytical access, security, and observability were assessed through retained runtime evidence.

Chapter IX builds on these results to examine the architectural strengths, trade-offs, limitations, and industrial relevance of the implemented solution.
